\scanpage{114}
\uline{Suppose} now that we have $\sum\alpha_i<+\infty$,
put $f(z)=\prod_1^\infty(1+\alpha_i z)$. Then $\rho$ is also
the order of the entire function, i.e
\[ \rho=\limsup_{0<r\to\infty}\frac{\log(\log f(r))}{\log r} \]

Indeed, we know that $\rho$ is $\leq$ order of the
function $f(z)$, \deleted{\unsure{and}} \unsure{let us show|let us prove} the reverse. We are
\uline{reduced} to showing that \unsure{.|:} \deleted{the product} in the
case where we have $\alpha_n=\frac1{n^\beta}$, we have
$\displaystyle \limsup_{0<r\to\infty}\frac{\log\log f(r)}{\log r}\leq\frac1\beta$ \quad (in short we in fact have $=\frac1\beta$)

Indeed, we have
\[ f(r)=\prod_1^\nu\left(1+\frac1{n^\beta}r\right)\prod_{\nu+1}^{\infty}\left(1+\frac1{n^\beta}r\right) \]
the second factor is $\leq e^{R(\nu)r}$
with $\displaystyle R(\nu)=\sum_{\nu+1}^\infty\frac1{n^\beta}\leq\int_\nu^\infty\frac{du}{x^\beta}=\frac1{\beta-1}\frac1{\nu^{\beta-1}}$
Take $\nu$ such that $r\unsure{\cdot}\frac1{\nu^{\beta-1}}\leq r^{\frac1\beta}$, i.e $\nu\geq r^{\frac1\beta}$
to be precise, $\nu=[r^{\frac1\beta}]+1$. Then, $e^{R(\nu)r}$
will be a \uline{function of $r$} of order $\leq\frac1\beta$\unsure{;|.} It
suffices to verify it for the product $\prod_1^\nu(1+\frac1{n^\beta}r)$
$\leq(1+r)^\nu=e^{\nu\log(1+r)}$, which is manifestly
of order $\leq\frac1\beta$.

Recall that if we put $f(\unsure{r|z})=\sum a_n z^n$,
\deleted{\unsure{in|therefore}} and $\unsure{B|D}(r)=\mathop{\rm Sup}_n(a_n r^n)$, then we
also have \qquad order \unsure{a|of} $f=\limsup_{r\to\infty}\frac{\log\log B(r)}{\log r}$
\margin{115}
